The system is built to support the evaluation described here, but that evaluation has \emph{not
been carried out}. This section is written in the future tense throughout and states a protocol,
not results. It is included because the design decisions in Section~\ref{sec:proposed_system} were
made against these measurements, and because specifying them makes the boundary between what has
been demonstrated and what remains unvalidated explicit.

\subsection{Design}

The evaluation separates into the four studies listed in Table~\ref{tab:studies}, which have different requirements. Combining them into a
single protocol, as is common, would force sensor characterisation and behavioural observation into
the same session even though the first needs no participants and the second needs unscripted work.

\begin{table}[htbp]
\centering
\caption{Planned studies and what each answers}
\label{tab:studies}
\begin{tabular}{@{}llp{6.0cm}@{}}
\toprule
\textbf{Study} & \textbf{Participants} & \textbf{Answers} \\
\midrule
A --- Bench validation      & none              & Sensor accuracy and angular repeatability \\
B --- Scripted conditions   & 8--10             & Detection performance, condition discrimination \\
C --- Free use              & 6--8, 2 sessions  & Recommendation quality, correction behaviour \\
D --- Replay ablation       & none (uses B, C logs) & Contribution of temporal qualification \\
\bottomrule
\end{tabular}
\end{table}

\subsection{Study A: Bench Validation}

Distance readings from both channels will be compared against a tape measure at fixed marks between
$30$ and $90\,\text{cm}$, reporting mean absolute error and root-mean-square error per channel and
the agreement between channels. Illuminance will be compared against a calibrated reference meter
across a range of lighting conditions; if no reference instrument is available, the claim reduces
to repeatability and to agreement with intended lighting categories, which is weaker but honest.

Head orientation requires a different treatment. Because angles are expressed relative to a
user-set zero (Section~\ref{subsec:vision}), there is no anatomical true value to be accurate
against, and a mean absolute error against an anatomical reference would be meaningless. The
measurable properties are \textbf{repeatability} --- the spread of reported angle across repeated
returns to the same marked head position --- and \textbf{sensitivity} --- whether a known angular
displacement $\Delta$ produces a reported change of $\Delta$. Both will be assessed against marked
positions on a fixed rig.

\subsection{Study B: Scripted Conditions}

Participants will hold each of the six conditions in Table~\ref{tab:conditions} for a fixed interval.

\begin{table}[htbp]
\centering
\caption{Experimental conditions}
\label{tab:conditions}
\begin{tabular}{@{}llp{6.6cm}@{}}
\toprule
\textbf{Condition} & \textbf{Scenario} & \textbf{Description} \\
\midrule
1 & Normal ergonomic condition & Appropriate posture, suitable distance, normal lighting \\
2 & Close viewing distance     & Normal posture, intentionally reduced user-screen distance \\
3 & Forward head posture       & Forward head position with otherwise normal conditions \\
4 & Excessive neck flexion     & Head tilted downward beyond the neutral range \\
5 & Low lighting               & Normal posture and distance under inadequate lighting \\
6 & Combined risk condition    & Forward head, close distance, and poor lighting together \\
\bottomrule
\end{tabular}
\end{table}

Ground truth comes from the assigned condition and from manual annotation of recorded video, never
from the system's own output. Reported metrics will be per-flag precision, recall and F1-score with
a confusion matrix, and, for blink detection, precision and recall against frame-level annotation
of a ten-minute segment per participant.

Eight to ten participants is sufficient to populate a confusion matrix across six conditions with
readable per-cell counts. It is not sufficient to claim generalisation across body size, eyewear,
or skin tone, and that limitation will be stated rather than implied.

\subsection{Study C: Free Use}

Participants will work normally for a fixed period with the system active. Two measures follow.

Recommendation quality will be rated on five-point Likert scales for relevance, clarity,
usefulness, and actionability, collected on a form and joined to the session record by session
identifier. To test whether the language model earns its dependency, the generative and
deterministic paths will be alternated within participant and counterbalanced, holding the
detection layer identical. A within-subject design is used deliberately: at these sample sizes a
between-subjects comparison has no realistic power.

Correction behaviour will be measured by the corrected-per-reminder ratio and time-to-correction
per flag type, both already instrumented. The unit of analysis here is the \textbf{reminder, not
the participant}: six to eight participants over two sessions each yields on the order of a hundred
reminders, which is a usable denominator even though the participant count is small. Results will
be reported descriptively; no claim of statistical significance will be made at this sample size.

\subsection{Study D: Replay Ablation}

Logs from Studies B and C will be replayed with the qualification layer disabled, isolating its
contribution to interruption load on identical input. This requires no additional data collection
and provides the comparison against a single-factor threshold baseline.

\subsection{Ethics}

The study involves continuous facial video and therefore requires informed consent covering: that
frames are transmitted off the device for processing, how long recordings are retained and who can
access them, that spoken audio prompts form part of the intervention, and the participant's right
to end a session at any point. Since the system is a screening and behavioural-support tool rather
than a diagnostic instrument, participants will be told explicitly that its output is not a health
assessment.
